\begin{lemma}[The operator domain sits inside the form domain: $\operatorname{Dom}(L^{1/2}) \subseteq \operatorname{Dom}(L^{1/4})$]
Let $v \in \operatorname{Dom}(L^{1/2})$ in the sense of \noderef{FracPowDomain}.
Then the fractional energy of \noderef{Energy} is finite,
\[
  \mathcal{E}(v) = \sum_{\alpha \neq 0} |\alpha|^{1/2}\,|v_\alpha|^2 < \infty ,
\]
where $v_\alpha$ are the Hermite coefficients of \noderef{HermiteBasis} and
$|\alpha|$ is the weight of \noderef{MultiIndexSum}.  Equivalently, $v$ belongs
to the form domain
$\mathcal{V} = \{w \in L^2(\gamma_n) : \mathcal{E}(w) < \infty\}
 = \operatorname{Dom}(L^{1/4})$ of \noderef{Energy}: the operator domain of
$B = L^{1/2}$ is contained in the form domain of the energy.  Only this
inclusion is asserted; the reverse inclusion is false and is not claimed.
\end{lemma}
\begin{proof}
Write $v_\alpha$ for the Hermite coefficients of $v$ relative to
\noderef{HermiteBasis}.

\emph{Step 1: what the hypothesis says.}
By the criterion of \noderef{FracPowDomain} at $s = 1/2$, whose multiplier is
$|\alpha|^{2s} = |\alpha|$, the hypothesis $v \in \operatorname{Dom}(L^{1/2})$
says exactly
\begin{equation*}
  \sum_{\alpha \neq 0} |\alpha|\,|v_\alpha|^2 < \infty . \tag{1}
\end{equation*}
All summands are nonnegative, since $|\alpha|$ is a nonnegative integer by
\noderef{MultiIndexSum}, so by the extended-real bridge clause of the
summable-families item of the shared context, $(1)$ is equivalent to
summability of the family $\bigl(|\alpha|\,|v_\alpha|^2\bigr)_{\alpha \neq 0}$
as a family of real numbers.

\emph{Step 2: $|\alpha| \geq 1$ for $\alpha \neq 0$.}
By \noderef{MultiIndexSum}, $|\alpha| = \sum_{j=1}^n \alpha_j$ with every
$\alpha_j \in \mathbb{N}_0$.  If $\alpha \neq 0$ then $\alpha_{j_0} \geq 1$ for
some $j_0$, and all the remaining summands are $\geq 0$, so
$|\alpha| \geq \alpha_{j_0} \geq 1$.

\emph{Step 3: the termwise comparison.}
Let $\alpha \neq 0$ and put $t = |\alpha| \geq 1$ by Step 2.  Then
$t^{1/2} \geq 1$, because $t^{1/2} < 1$ would give
$t = t^{1/2}\,t^{1/2} < 1$; and multiplying the inequality $1 \leq t^{1/2}$ by
the nonnegative number $t^{1/2}$ gives $t^{1/2} \leq t$.  Hence, multiplying by
$|v_\alpha|^2 \geq 0$,
\begin{equation*}
  0 \leq |\alpha|^{1/2}\,|v_\alpha|^2 \leq |\alpha|\,|v_\alpha|^2
  \qquad (\alpha \neq 0). \tag{2}
\end{equation*}

\emph{Step 4: conclusion.}
By Step 1 the majorant family in $(2)$ is summable, so by the comparison clause
of the summable-families item of the shared context the family
$\bigl(|\alpha|^{1/2}|v_\alpha|^2\bigr)_{\alpha \neq 0}$ is summable as well,
with
\[
  \sum_{\alpha \neq 0} |\alpha|^{1/2}\,|v_\alpha|^2
    \leq \sum_{\alpha \neq 0} |\alpha|\,|v_\alpha|^2 < \infty .
\]
Its summands being nonnegative, the extended-real bridge clause applies once
more and says that the same sum, read in $[0,+\infty]$, is finite and equal to
the real number just bounded.  By \noderef{Energy} that extended-real sum is
precisely $\mathcal{E}(v)$, so $\mathcal{E}(v) < \infty$.  By the description of
the form domain in \noderef{Energy} this is the assertion
$v \in \mathcal{V} = \operatorname{Dom}(L^{1/4})$.
\end{proof}
